\documentclass[UTF8,fontset=fandol,11pt]{ctexart}
\usepackage[a4paper,margin=25mm]{geometry}
\usepackage{amsmath,amssymb,amsthm,mathtools}
\usepackage{xcolor}
\usepackage{booktabs,array,enumitem,fancyhdr,hyperref}
\hypersetup{colorlinks=true,linkcolor=blue!50!black,urlcolor=blue!50!black}
\setlength{\emergencystretch}{2em}
\setlength{\parskip}{0.3em}
\linespread{1.16}
\allowdisplaybreaks
\newtheorem{theorem}{定理}[section]
\newtheorem{proposition}[theorem]{命题}
\newtheorem{lemma}[theorem]{引理}
\theoremstyle{definition}
\newtheorem{definition}[theorem]{定义}
\theoremstyle{remark}
\newtheorem{remark}[theorem]{说明}
\newcommand{\dd}{\,\mathrm d}
\newcommand{\ii}{\mathrm i}
\newcommand{\Rea}{\operatorname{Re}}
\newcommand{\Ima}{\operatorname{Im}}
\newcommand{\Li}{\operatorname{Li}}
\newcommand{\supp}{\operatorname{supp}}
\pagestyle{fancy}
\fancyhf{}
\fancyhead[L]{\small 学习笔记}
\fancyhead[R]{\small 初稿}
\fancyfoot[C]{\thepage}
\setlength{\headheight}{15pt}
\author{}
\date{2026年10月22日初稿}
\makeatletter
\renewcommand{\maketitle}{%
  \thispagestyle{plain}%
  \begin{center}
    {\LARGE\bfseries \@title\par}
    \vspace{0.6em}
    {\small \@date\par}
  \end{center}
  \vspace{0.6em}}
\makeatother
\newcommand{\draftnotice}{\noindent{\small
本文为供逐篇审阅的学习笔记初稿，不是新增研究成果。
文中区分直接证明与引用结果；正文不依赖本文。
本轮已作公式、条件及排版初步核查，尚待作者逐句审定。}\par\medskip}

\title{从 Euler 乘积到素数 $\zeta$ 函数}
\begin{document}
\maketitle
\draftnotice
\section{只对素数求和会得到什么}
$\zeta(s)$ 对全部正整数求和，而 Euler 乘积对素数取乘积。
若改为直接对素数求和，可以定义
\[
 P(s)=\sum_p p^{-s},\qquad \Rea s>1.
\]
这个函数通常称素数 zeta 函数。它与 $\log\zeta$ 密切相关，
但两者不是同一个函数；区分素数与素数幂，是理解这种关系的关键。
本篇在绝对收敛区域内完整推导两者的互反，并研究 $s\downarrow1$ 的实轴行为。
不声称据此已经完成复杂的全域延拓。

\section{由 Euler 乘积获得第一重关系}
在 $\sigma=\Rea s>1$ 上取与实轴正值连续一致的解析对数，
\[
 \log\zeta(s)=\sum_p\sum_{k=1}^\infty\frac{p^{-ks}}k
            =\sum_{k=1}^\infty\frac{P(ks)}k.                \tag{1}
\]
换序由绝对收敛保证：
\[
 \sum_p\sum_{k\ge1}\frac{p^{-k\sigma}}k
 \le\sum_p\frac{p^{-\sigma}}{1-p^{-\sigma}}
 \le\frac1{1-2^{-\sigma}}\sum_{n\ge2}n^{-\sigma}<\infty.
\]
$k=1$ 是素数项，$k\ge2$ 收集更高素数幂，故 $\log\zeta$ 并不等于 $P$。

\section{Möbius 反演的直接证明}
记 $\mu$ 为 Möbius 函数，其基本恒等式为
$\sum_{d\mid m}\mu(d)=1$（$m=1$），而 $m>1$ 时为零。
将式(1)代入，候选反演为
\[
 \sum_{k\ge1}\frac{\mu(k)}k\log\zeta(ks)
 =\sum_{k,j\ge1}\frac{\mu(k)}{kj}P(kjs).
\]
为了按 $m=kj$ 合并，先估计绝对值。对 $m\ge1$，
\[
 P(m\sigma)\le
 \sum_{n\ge2}n^{-m\sigma}
 \le 2^{-(m-1)\sigma}\sum_{n\ge2}n^{-\sigma}.
\]
因此
\[
 \sum_{k,j\ge1}\frac{|P(kjs)|}{kj}
 \le\sum_{m\ge1}\frac{d(m)}m P(m\sigma)<\infty,
\]
因为 $d(m)\le m$，而剩下的是几何衰减。
合法合并后，$P(ms)$ 的系数为
$m^{-1}\sum_{k\mid m}\mu(k)$，故
\[
 P(s)=\sum_{k=1}^\infty\frac{\mu(k)}k\log\zeta(ks),
 \qquad\Rea s>1.                                        \tag{2}
\]
这里不是把一个形式反演未经检查地用于双重无穷和。

\section{实轴上靠近一的行为}
对实数 $s>1$，将式(1)写成
\[
 P(s)=\log\zeta(s)-H(s),\qquad
 H(s)=\sum_p\sum_{k\ge2}\frac{p^{-ks}}k.
\]
当 $s\ge1$ 时，
\[
 0\le H(s)\le\sum_{n\ge2}\sum_{k\ge2}n^{-k}
 =\sum_{n\ge2}\frac1{n(n-1)}=1.
\]
控制收敛定理因而保证 $H(s)\to H(1)<\infty$。
又由 $\zeta$ 在 $1$ 的留数为 $1$，
$\zeta(s)=(s-1)^{-1}(1+O(s-1))$，所以
\[
 P(s)=\log\frac1{s-1}-H(1)+o(1),\qquad s\downarrow1.
\]
$\zeta$ 的简单极点经取对数后变成对数型发散；高次素数幂只改变有限常数。
这已经说明两个对象的奇性类型不同，但不是复平面上的完整分支分析。

\section{从收敛边界能推出什么，不能推出什么}
若假设 $\sum_p1/p<\infty$，则对 $s\downarrow1$ 用控制收敛定理，
$P(s)$ 应趋于有限值，与上节发散矛盾。所以
$\sum_p1/p$ 发散。
更弱的“素数有无穷多个”也随之成立。

另一方面，式(2)中出现 $\log\zeta(ks)$，
提醒我们继续延拓时必须处理零点、极点与对数分支。
但不能仅看见一个候选奇点，就宣称它一定在最终的和中保留下来：
不同项之间可能存在重合或抵消，还须证明。
本篇因此只在已建立绝对收敛的区域使用式(2)，不把后续奇性位置画成未经论证的结论。

\section*{参考依据与后续核验}
Euler 乘积、$\zeta$ 在 $1$ 的极点及 Möbius 基本恒等式，可参看本书相关结果。
本篇直接证明了双重和反演和实轴渐近。
若以后讨论自然边界，必须另列所用的零点信息和分支构造，不能作为本篇末尾的一句附带断言。
\end{document}
